\section{Experiments}
The two different backbone architecture for the generation of the latent space representation of the objects and the performance of the different clustering algorithms are evaluated by hyperparameter optimization of the different parameters along the entire pipeline of the procedure is evaluated in this section.

\subsection{Baseline algorithm}
The already existing PointNet autoencoder currently used at \ac{IPA} is used as a baseline for evaluating the performance of the proposed algorithms in this thesis. The dimension of the output of the latent space representation of the baseline algorithm is 64. As mentioned in the restrictions in Section. \ref*{sec:restrictions}, usually a two-step dimensionality reduction is performed before the datasets are grouped into clusters for the two networks evaluated in this thesis because the dimension of the output of the latent space representations is too high. This is done to suppress some noise and also to escalate the procedure by faster calculation of the pairwise distance between the datapoints. But since an output dimension of 64 is not very high, two-step dimensionality reduction is not necessary. However, to make the baseline algorithm comparable with the proposed networks, one-step dimensionality reduction is performed in some cases. 

\subsubsection{Basic Study without dimensionality reduction}
Table \ref*{tab:basic_without_tsne} shows the results of the baseline algorithm when no dimensionality reduction technique is used. It shows the results of the basic algorithm for the clustering of 1000 datapoints. The unit of time across all the tables in this section is in seconds unless otherwise mentioned. The four clustering algorithms \ac{HDBSCAN}, Kmeans clustering, spectral clustering and agglomerative clustering are evaluated. The number of clusters generated by the \ac{HDBSCAN} algorithm cannot be fixed by the user and hence it is determined by the algorithm automatically. Thus, the number of clusters after the first level of the multi-level \ac{HDBSCAN} algorithm as proposed in Section. \ref{sec:clustering_algo} is 232 in this case. As mentioned in Section \ref{sec:restrictions} the number of clusters for the rest of the clustering algorithms is set to 175 throughout all the experiments in this section. The \ac{DBCV} score caries in the range [-1,1] with more positive results denoting better clustering of the datapoints. The reason why \ac{HDBSCAN} algorithm shoes better results as compared to other clustering algorithms is \ac{HDBSCAN} akgorithm is a density based clustering algorithm. \ac{DBCV} being a density based clustering validation index is likely to give better results for density based clustering algorithms which is proven to be true in this experiment. Furthermore, theoretically \ac{HDBSCAN} is better capable of handling noises and outliers in the dataset which can also contribute to the better results as exhibited in this experiment. \todo{include experiment on multi-level HDBSCAN}
\begin{table}[ht]
    \setlength\extrarowheight{10pt}
    \caption{Results of the baseline algorithm on 1000 datapoints and without any dimensionality reduction technique. The unit of time is in seconds.}
    \centering
    \begin{tabular}{l|l|l|l|l|l}
      \toprule
      No. of datapoints & Clustering Algo. & No. of Clusters & Clustering time & \ac{DBCV} score & \ac{DBCV} time\\
      \midrule
      \multirow{4}{1.0in}{1000} & HDBSCAN & 232 & 0.182 & -0.189 & 297.031 \\ \cline{2-6} 
                                & Kmeans & 175	& 3.960 & -0.332 & 237.640 \\ \cline{2-6} 
                                & Spectral & 175 & 2.396 & -0.339	& 161.604 \\ \cline{2-6}
                                & Agglomerative	& 175 & 0.115 & -0.327 & 157.584  \\ 
      \bottomrule
    \end{tabular}
    \label{tab:basic_without_tsne}
\end{table}

\subsubsection{Basic Study with \ac{t-SNE}}
Table \ref*{tab:basic_with_tsne} shows the results of the baseline algorithm when the \ac{t-SNE} dimensionality reduction is applied. The number of components for \ac{t-SNE} could be 1, 2 or 3 as per the available scikit-learn implementation \cite*{scikit-learn}. It was chosen to be 2 in these set of experimentation for better and easier visualization. The \ac{DBCV} score caries in the range [-1,1] with more positive results denoting better clustering of the datapoints. The results are seen to be significantly better as compared to Table \ref*{tab:basic_without_tsne} when dimensionality reduction technique was not applied. This is because the \ac{t-SNE} algoithm tends to keep datapoints that are distant from one another in the higher dimensional space also distant datapoints even in the lower-dimensional space. This gives rise to more dense and distinct clusters thus improving the \ac{DBCV} score. As opposed to the result in Tab.\ref*{tab:basic_without_tsne}, where the \ac{DBCV} score was negative, meaning the datapoints were mostly wrongly classification, in this table it can be seen that the datapoints are mostly correctly classified. It is also to be note that the number of clusters generated by the \ac{HDBSCAN} algorithm is 286 which is a bit more as compared to the previous case. This further hints to the direction that as a result of using \ac{t-SNE}, more dense and non-overlapping clusters are formed in the dataset. 
\begin{table}[ht]
    \setlength\extrarowheight{10pt}
    \caption{Results of the baseline algorithm on 1000 datapoints after applying \ac{t-SNE} dimensionality reduction. The unit of time is in seconds.}
    \centering
    \begin{tabular}{p{30pt}|l|l|l|p{50pt}|p{40pt}|p{30pt}|p{30pt}}
      \toprule
      No. of datapoints & \ac{t-SNE} dim	& \ac{t-SNE} time & Clustering Algo. & No. of Clusters & Clustering time & \ac{DBCV} score & \ac{DBCV} time\\
      \midrule
      \multirow{4}{30pt}{1000} & 2	& 3.030 & HDBSCAN	& 286	& 0.110 & 0.0975	& 275.039 \\ \cline{2-8} 
                                & 2	& 2.777 &	Kmeans	& 175	& 0.191 & -0.052  	& 94.004 \\ \cline{2-8} 
                                & 2	& 3.318 & Spectral	& 175	& 3.604 & -0.064 	& 126.854 \\ \cline{2-8}
                                & 2	& 3.190 & Agglomerative	& 175	& 0.087 & -0.076	& 113.937 \\ 
      \bottomrule
    \end{tabular}
    \label{tab:basic_with_tsne}
\end{table}
\subsubsection{Basic Study with increasing datapoints}
As mentioned in Section \ref*{sec:restrictions}, it is extremely inefficient to compute the \ac{DBCV} for all the datapoints in the ABC datasets \cite*{Koch_2019_CVPR} as once, it was necessary to verify that the results obtained for a smaller subset of the dataset can be interpolated for the entire dataset. In other words, the results obtained for a smaller subset of the dataset generalises the behavior of the entire dataset well. Table \ref*{tab:basic_2000_without_tsne} shows the results of the baseline algorithm on 2000 datapoints without applying dimensionality reduction and table \ref*{tab:basic_2000_with_tsne} shows the results of the baseline algorithm on 2000 datapoints when \ac{t-SNE} dimensionality reduction is applied. It was observed that roughly $10\%$ of the datapoints were classified as outliers by the \ac{HDBSCAN} algorithm for all the experiments on 1000 and 2000 datapoints. 
\begin{table}[ht]
    \setlength\extrarowheight{10pt}
    \caption{Results of the baseline algorithm on 2000 datapoints and without any dimensionality reduction technique. The unit of time is in seconds.}
    \centering
    \begin{tabular}{l|l|l|l|l|l}
      \toprule
      No. of datapoints & Clustering Algo. & No. of Clusters & Clustering time & \ac{DBCV} score & \ac{DBCV} time\\
      \midrule
      \multirow{4}{1.0in}{2000} & HDBSCAN	& 436	& 0.444 & -0.234	& 2115.777 \\ \cline{2-6} 
                                & Kmeans	& 175	& 0.793 & -0.462	& 764.222 \\ \cline{2-6} 
                                & Spectral	& 175	& 20.607 & -0.617	& 749.487 \\ \cline{2-6}
                                & Agglomerative	& 175	& 0.205 & -0.502	& 756.480 \\ 
      \bottomrule
    \end{tabular}
    \label{tab:basic_2000_without_tsne}
\end{table}

\begin{table}[ht]
    \setlength\extrarowheight{10pt}
    \caption{Results of the baseline algorithm on 2000 datapoints after applying \ac{t-SNE} dimensionality reduction. The unit of time is in seconds.}
    \centering
    \begin{tabular}{p{30pt}|l|l|l|p{50pt}|p{40pt}|p{30pt}|p{30pt}}
      \toprule
      No. of datapoints & \ac{t-SNE} dim	& \ac{t-SNE} time & Clustering Algo. & No. of Clusters & Clustering time & \ac{DBCV} score & \ac{DBCV} time\\
      \midrule
      \multirow{4}{30pt}{2000} & 2	& 6.689 & HDBSCAN	& 626	& 0.237 secs	& 0.172	& 4497.151 \\ \cline{2-8} 
                                & 2	& 7.450 & Kmeans	& 175	& 0.155 & -0.226 &	932.509 \\ \cline{2-8} 
                                & 2	& 157.017& Spectral	& 175	& 91.684 & -0.125 & 755.647 \\ \cline{2-8}
                                & 2	&  & Agglomerative	\\ 
      \bottomrule
    \end{tabular}
    \label{tab:basic_2000_with_tsne}
\end{table}

\subsection{Conclu approach with PointNet as Backbone}
As mentioned in Section \ref*{sec:method}, wo feature representation learning backbones were used for the experiments. Experiments with the PointNet autoencoder as the backbone is recorded here. The performanceof all the clustering algorithm has been evaluated. Different hyperparameter optimization has been performed which will be elaborated  gradually through this subsection.  

\subsubsection{Evaluation of increasing datapoints}
As mentioned in Section \ref*{sec:restrictions}, it is extremely inefficient to compute the \ac{DBCV} for all the datapoints in the ABC datasets \cite*{Koch_2019_CVPR} as once, it was necessary to verify that the results obtained for a smaller subset of the dataset can be interpolated for the entire dataset. In other words, the results obtained for a smaller subset of the dataset generalises the behavior of the entire dataset well. Table \ref*{tab:conclu_1000} shows the results of the ConClu approach on 1000 datapoints and table \ref*{tab:conclu_2000}  and \ref*{tab:conclu_5000} shows the results of the ConClu approach on 2000 datapoints and 5000 datapoints respectively. It was observed that roughly $10\%$ of the datapoints were classified as outliers by the \ac{HDBSCAN} algorithm for all the experiments on 1000, 2000 and 5000 datapoints. The dimension of the latent space representations for table
\ref*{tab:conclu_1000}, \ref*{tab:conclu_2000}, and \ref*{tab:conclu_5000} is 1024. 
\begin{table}[ht]
    \setlength\extrarowheight{10pt}
    \caption{Results of the ConClu approach on 1000 datapoints. The unit of time is in seconds.}
    \centering
    \begin{tabular}{l|l|l|l|l|l}
      \toprule
      No. of datapoints & Clustering Algo. & No. of Clusters & Clustering time & \ac{DBCV} score & \ac{DBCV} time\\
      \midrule
      \multirow{4}{1.0in}{1000} & HDBSCAN	& 288	& 0.228 & 0.137	& 437.956 \\ \cline{2-6} 
                                & Agglomerative	& 175	& 0.063 & -0.057	& 191.887 \\ \cline{2-6} 
                                & Kmeans	& 175	& 7.763 & -0.120 	& 192.580 \\ \cline{2-6}
                                & Spectral	& 175	& 28.831 & -0.051	& 186.879 \\ 
      \bottomrule
    \end{tabular}
    \label{tab:conclu_1000}
\end{table}

It can be noticed that the \ac{DBCV} score for the \ac{HDBSCAN} algorithm increases consistently from 1000 datapoints to 2000 datapoints and so on. This could be attributed to the fact that more datapoints would quite naturally give rise to denser clusters and more dense the inside of the cluster is as compared to its outside, the better is the \ac{DBCV} score. This further reassures the fact that denser clusters would mean more number of elements in individual clusters. In other words, the representative member of the cluster would represent a more diverse amount of other cluster members.  

\begin{table}[ht]
    \setlength\extrarowheight{10pt}
    \caption{Results of the ConClu approach on 2000 datapoints. The unit of time is in seconds.}
    \centering
    \begin{tabular}{l|l|l|l|l|l}
      \toprule
      No. of datapoints & Clustering Algo. & No. of Clusters & Clustering time & \ac{DBCV} score & \ac{DBCV} time\\
      \midrule
      \multirow{4}{1.0in}{2000} & HDBSCAN	& 605	& 1.030 & 0.159	& 8278.394 \\ \cline{2-6} 
                                & Agglomerative	& 175	& 0.196 & -0.212	& 786.866 \\ \cline{2-6} 
                                & Kmeans	& 175	& 6.322 & -0.209	& 1079.195 \\ \cline{2-6}
                                & Spectral	& 175	& 80.008 & -0.145	& 1187.477 \\ 
      \bottomrule
    \end{tabular}
    \label{tab:conclu_2000}
\end{table}

Also it is noteworthy how the time to calculate the \ac{DBCV} score doesn't increase linearly on increasing the number of datapoints. Rather it increases in an exponential rate. This shows how it would take immensely long to compute the \ac{DBCV} score of the entire dataset in the presence of 1 million datapoints in it. But since the results obtained on a smaller subset of the dataset can be interpolated to the entire dataset, evaluation hyperparameter optimization for the rest of the parameters are performed on 1000 datapoints. However, after deciding the final values for all the hyperparameters, the cluster representatives are generated taking into account all the datapoints under consideration as that no longer involves calculating the \ac{DBCV} score anymore for evaluation.  
\begin{table}[ht]
    \setlength\extrarowheight{10pt}
    \caption{Results of the ConClu approach on 5000 datapoints. The unit of time is in seconds.}
    \centering
    \begin{tabular}{l|l|l|l|l|l}
      \toprule
      No. of datapoints & Clustering Algo. & No. of Clusters & Clustering time & \ac{DBCV} score & \ac{DBCV} time\\
      \midrule
      \multirow{4}{1.0in}{5000} & HDBSCAN	& 1.586	& 1.904 & 0.210	& 165884.854 \\ \cline{2-6} 
                                & Agglomerative	& 175	& 1.345 & -0.380	& 4945.410 \\ \cline{2-6} 
                                & Kmeans	& 175	& 9.142 & -0.427	& 8289.148 \\ \cline{2-6}
                                & Spectral	& 175	& 641.189 & -0.359	& 5069.388 \\ 
      \bottomrule
    \end{tabular}
    \label{tab:conclu_5000}
\end{table}

\subsubsection{Evaluation of different dimensionality reduction techniques}
As mentioned in section \ref*{sec:restrictions}, it is essential to perform dimensionality reduction before clustering the objects in the dataset into groups. Two different dimensionality techniques were evaluated for this purpose, the \ac{PCA} method and the \ac{t-SNE} method. Table \ref*{tab:conclu_pca_1000} shows the results on the ConClu approach for 1000 datapoints and applying \ac{PCA} method as the dimensionality reduction technique in both steps. As the \ac{DBCV} score depends on the number of nearest neighbor raised to the power of the negative of 1 divided by the number of features in the latent space representation as shown in Eq. \ref*{eq:aptscoredist}, it is important to determine the to the maximal possible dimension of the representation after the second step of dimensionality reduction. Experiments were performed and the value was found to be 154 in this use case. Also keeping the first step of dimensionality reduction to be \ac{PCA}, experiments were also conducted with \ac{t-SNE} method as the second step. Table \ref*{tab:conclu_tsne_1000} shows the results on the ConClu approach for 1000 datapoints under such settings. The dimensionof the latent space representations in both the cases is set to 1024. it is to be noted that the results for all the clustering algorithms have better performance when \ac{t-SNE} is used as the dimensionality reduction technique in the second step as compared to \ac{PCA}. It can be attributed to the fact that the \ac{t-SNE} algorithm preserves the spatial relationship between the datapoints. In other words, if datapoints are close to one another in a higher dimensional space, \ac{t-SNE} tends to preserve this spatial closeness between points even in the lower dimensional space. 

\begin{table}[ht]
    \setlength\extrarowheight{10pt}
    \caption{Results of the Conclu approach with PointNet autonocoder on 1000 datapoints with \ac{PCA} as dimensionality reduction technique in both steps. The unit of time is in seconds.}
    \centering
    \begin{tabular}{p{30pt}|p{30pt}|p{30pt}|p{30pt}|l|p{30pt}|p{40pt}|p{30pt}|p{30pt}}
      \toprule
      \ac{PCA} dim & \ac{PCA} time & \ac{PCA} dim	& \ac{PCA} time & Clustering Algo. & No. of Clusters & Clustering time & \ac{DBCV} score & \ac{DBCV} time\\
      \midrule
      \multirow{4}{30pt}{512}	& 32.921 & 154	& 11.249 & HDBSCAN	& 169	& 0.459 & -0.410	& 234.172 \\ \cline{2-9} 
                                & 41.230 & 154	& 14.891 & Agglomerative	& 175	& 0.216 & -0.151	& 209.954 \\ \cline{2-9} 
                                & 38.408 & 154	& 13.032 & Kmeans	& 175	& 7.001 & -0.184	& 212.482 \\ \cline{2-9}
                                & 44.171 & 154	& 13.897 & Spectral	& 175	& 46.703 &-0.481 & 221.064 \\ 
      \bottomrule
    \end{tabular}
    \label{tab:conclu_pca_1000}
\end{table}

Also it is to be noted that the \ac{HDBSCAN} algorithm performs significantly bad in presence of \ac{PCA} as compared to \ac{t-SNE}. This is also because of the property of \ac{t-SNE} mentioned above which gives rise to denser clusters as compared to \ac{PCA}.  

\begin{table}[ht]
    \setlength\extrarowheight{10pt}
    \caption{Results of the Conclu approach with PointNet autonocoder on 1000 datapoints with \ac{PCA} as dimensionality reduction technique in the first step and \ac{t-SNE} as the second. The unit of time is in seconds.}
    \centering
    \begin{tabular}{p{30pt}|p{30pt}|p{30pt}|p{30pt}|l|p{30pt}|p{40pt}|p{30pt}|p{30pt}}
      \toprule
      \ac{PCA} dim & \ac{PCA} time & \ac{t-SNE} dim	& \ac{t-SNE} time & Clustering Algo. & No. of Clusters & Clustering time & \ac{DBCV} score & \ac{DBCV} time\\
      \midrule
      \multirow{4}{30pt}{512}	& 32.530 & 2	& 143.024 & HDBSCAN	& 288	& 0.228 & 0.137	& 437.956 \\ \cline{2-9} 
                                & 38.832 & 2	& 141.566 & Agglomerative & 175	& 0.063 & -0.057 & 191.887 \\ \cline{2-9} 
                                & 39.326 & 2	& 162.957 & Kmeans	& 175	& 7.763 & -0.120 	& 192.580 \\ \cline{2-9}
                                & 71.552 & 2	& 177.427 & Spectral	& 175	& 28.831 & -0.051	& 186.879 \\ 
      \bottomrule
    \end{tabular}
    \label{tab:conclu_tsne_1000}
\end{table}

The effect of two step dimensionality reduction was also compared against only one step dimensionality reduction as shown in table \ref*{tab:conclu_1000_without_reduction}. A \ac{PCA} followed by a \ac{t-SNE} shows better results as comoared to directly applying \ac{t-SNE} just once. This is because of the fact that a multi-step dimensionality reduction helps in reducing the noises during the procedure.

\begin{table}[ht]
    \setlength\extrarowheight{10pt}
    \caption{Results of the Conclu approach with PointNet autonocoder on 1000 datapoints with only \ac{t-SNE} as dimensionality reduction technique as the only  step. The unit of time is in seconds.}
    \centering
    \begin{tabular}{p{30pt}|l|l|l|p{50pt}|p{40pt}|p{30pt}|p{30pt}}
        \toprule
        No. of datapoints & \ac{t-SNE} dim	& \ac{t-SNE} time & Clustering Algo. & No. of Clusters & Clustering time & \ac{DBCV} score & \ac{DBCV} time\\
        \midrule
        \multirow{4}{30pt}{1000} & 2	& 6.769 & HDBSCAN	& 314	& 0.189 & -0.120 & 337.313 \\ \cline{2-8} 
                                  & 2	& 4.872 & Kmeans	& 175	& 0.337 & -0.132
                                  & 131.622 \\ \cline{2-8} 
                                  & 2	& 5.831 secs	& Spectral	& 175	& 7.592 & -0.122	& 130.678 \\ \cline{2-8}
                                  & 2	& 4.478 secs	& Agglomerative	& 175	& 0.033 & -0.085 & 128.358 \\ 
        \bottomrule
    \end{tabular}
    \label{tab:conclu_1000_without_reduction}
\end{table}

\subsubsection{Evaluation of dimension of feature space representations}
Latent space representations are lower dimensional representations of the dataset. It is expensive to work with very high dimensional data. Thus the dimension of the latent space representations is immensely important. It is extremely necessary that the original high dimensional data is compressed to a lower dimensional latent space without severe loss of information in the process. Since the original dataset had the point clouds with 2048, for the sake of completeness of the experimentation, the evaluation has been performed for all dimensions on a logarithmic scale of 2 from 2048 to down until 4. Table \ref*{tab:conclu_2048} shows the experiments when the dimension of the latent space representations was set to be 2048. As mentioned in Section \ref*{sec:restrictions}, because of the very high number of features in the output of the network, a two-step dimensionality reduction was performed to suppress noise during the process of dimensionality reduction and also to facilitate faster calculation of the pairwise distance between the points. 
\begin{table}[ht]
    \setlength\extrarowheight{10pt}
    \caption{Results of the Conclu approach with PointNet autonocoder on 1000 datapoints when the dimension of the latent space representations is 2048. The unit of time is in seconds.}
    \centering
    \begin{tabular}{p{30pt}|p{30pt}|p{30pt}|p{30pt}|l|p{30pt}|p{40pt}|p{30pt}|p{30pt}}
      \toprule
      \ac{PCA} dim & \ac{PCA} time & \ac{t-SNE} dim	& \ac{t-SNE} time & Clustering Algo. & No. of Clusters & Clustering time & \ac{DBCV} score & \ac{DBCV} time\\
      \midrule
      512	& 18.507 & 2	& 3.594 & HDBSCAN	& 289	& 0.127 & 0.189 & 243.297 \\ \cline{1-9} 
      512	& 15.525 & 2	& 3.368 & Agglomerative	& 175	& 0.036	& -0.013	& 109.668 \\ \cline{1-9} 
      512	& 16.595 & 2	& 3.709 & Kmeans & 175	& 0.222 & -0.040	& 102.518 \\ \cline{1-9}
      512	& 18.593 & 2	& 3.918 & Spectral	& 175	& 9.257 & 0.020	& 101.925 \\ 
      \bottomrule
    \end{tabular}
    \label{tab:conclu_2048}
\end{table}